% Generated draft: requires mathematical and source review.
\documentclass[11pt]{article}
\usepackage[margin=25mm]{geometry}
\usepackage{fontspec}
\setmainfont{Latin Modern Roman}
\usepackage{amsmath,amssymb,mathtools,unicode-math}
\setmathfont{Latin Modern Math}
\usepackage{enumitem,needspace,xurl,hyperref}
\setlength{\parindent}{0pt}
\setlength{\parskip}{5pt}
\setlength{\emergencystretch}{3em}
\newcommand{\N}{\mathbb{N}}
\newcommand{\Z}{\mathbb{Z}}
\newcommand{\Q}{\mathbb{Q}}
\newcommand{\R}{\mathbb{R}}
\newcommand{\C}{\mathbb{C}}
\newcommand{\F}{\mathbb{F}}
\newcommand{\A}{\mathbb{A}}
\newcommand{\PP}{\mathbb{P}}
\newcommand{\E}{\mathbb{E}}
\newcommand{\eps}{\varepsilon}
\newcommand{\dd}{\,\mathrm{d}}
\begin{document}
\section*{UnsolvedMath ID: 3700019}
\textbf{Source problem code:} catalogue entry
\subsection{Definitions and mathematical statement}
Let $F:\C^2\to\C$ be entire, let $z_0,w_0\in\C$, and let $\phi$ be a holomorphic germ at $z_0$ such that
\[
 \phi(z_0)=w_0,\qquad F(z,\phi(z))=0
\]
near $z_0$. The source assumes a nonsingular base point. For the usual regular implicit-function interpretation this means
\[
 F(z_0,w_0)=0,\qquad \frac{\partial F}{\partial w}(z_0,w_0)\ne0.
\]
The source does not explicitly distinguish this condition from nonsingularity of the curve $F=0$; the intended regularity convention therefore requires confirmation.

For $\theta\in[0,2\pi)$, write $r_\theta(t)=z_0+t e^{i\theta}$, $t\ge0$. Analytic continuation of the germ along the whole ray means continuation along $r_\theta|_{[0,T]}$ for every finite $T>0$, with the continuations agreeing along the ray. No value or regularity at infinity is required. The Gross-property question is whether every such implicit germ satisfies
\[
 \lambda\bigl(\{\theta\in[0,2\pi):
 \phi\text{ cannot be analytically continued along }r_\theta\}\bigr)=0,
\]
where $\lambda$ is Lebesgue angular measure. Equivalently, continuation should be possible along almost every ray issuing from its base point.

The supplied assessment distinguishes this assertion from the weaker Iversen property. In the latter, every curve $\gamma$ from $z_0$ can, for each $\varepsilon>0$, be replaced by a nearby path $\gamma_1$ along which $\phi$ continues, with the source writing $|\gamma-\gamma_1|\le\varepsilon$. The curve class, common parametrization, and precise distance convention for that auxiliary property are not supplied. This approximation property does not replace the almost-every-ray target.
\subsection{Short English statement}
Must a local holomorphic solution of an entire equation in two complex variables continue along almost every ray from its initial point? The supplied summary reports that the Gross property is known for inverse germs, corresponding to $F(z,w)=z-f(w)$ with $f$ entire, while Stöilov proved the weaker Iversen path-approximation property for implicit functions. It describes the general two-variable Gross assertion as open in the literature it checked and labels the record partial progress. These are supplied assessments, not new literature verification.

Provenance says the problem was retained after applying explicit solved/counterexample updates in the author's index through the access date 2026-07-29. That procedure and the registry label do not establish that no later or unindexed resolution exists.
\subsection{Sources}
\begin{itemize}
\item Local UnsolvedMath registry, permanent ID 3700019, source code catalogue entry, ``Gross property of implicit functions,'' attributed to Alexandre Eremenko and cited proposers, proposed year recorded as 2015. Dataset curation and metadata: Ulam AI and UnsolvedMath Contributors, CC BY 4.0.
\item Eremenko, \emph{My favorite unsolved problems}, item ``Gross property of implicit functions,'' \url{https://www.math.purdue.edu/~eremenko/uns1.html}; accessed 2026-07-29, extraction recorded as pdf-text.
\item Eremenko, ``Singularities of implicit functions,'' cited in the supplied summary by the filename \texttt{gross2.pdf}; no complete item-specific URL or publication details are supplied.
\item The summary attributes the Iversen property to Stöilov and also points to the companion on singularities of inverse functions, arXiv:2110.06134. No separate Stöilov bibliographic entry is supplied.
\item Source provenance marks status evidence NEEDS\_REVIEW and rights NEEDS\_REVIEW: public access was verified, but redistribution permission was not established. The source-specific qualification is separate from the dataset licence.
\end{itemize}
\end{document}
